\clearpage
\VerdierPhysicalPageAnchor{0269}%
\begingroup
\small
\setlength{\parindent}{0pt}%
\setlength{\parskip}{0pt}%
\setlength{\emergencystretch}{3em}%
\newcommand{\VerdierContentsRow}[3]{%
  \par\noindent
  \hangindent=#1\relax\hangafter=1
  \hspace*{#1}#2\nobreak
  \leaders\hbox{\kern.32em.\kern.32em}\hfill
  \nobreak\hbox{\kern.65em #3}\par}
\newcommand{\VerdierContentsChapter}[1]{%
  \par\addvspace{.55\baselineskip}\noindent\textbf{#1}\par\smallskip}
\newcommand{\VerdierContentsSection}[2]{%
  \par\addvspace{.45\baselineskip}\noindent
  \hfill\textbf{#1.}\enspace #2\hfill\mbox{}\par\smallskip}

\VerdierSemanticAnchor{toc-ch2}%
\VerdierContentsChapter{第II章：三角圏}
\VerdierSemanticAnchor{toc-ch2-sec1}%
\VerdierContentsSection{1}{定義と初歩的性質}
\VerdierSemanticAnchor{toc-ch2-sec1-1}%
\VerdierContentsRow{.5em}{\textbf{1.1.}\enspace 三角圏の定義}{93}
\VerdierSemanticAnchor{toc-ch2-sec1-2}%
\VerdierContentsRow{.5em}{\textbf{1.2.}\enspace 三角圏の初歩的性質。}{97}
\VerdierSemanticAnchor{toc-ch2-sec1-3}%
\VerdierContentsRow{.5em}{\textbf{1.3.}\enspace 三角圏の例。}{106}

\VerdierSemanticAnchor{toc-ch2-sec2}%
\VerdierContentsSection{2}{三角圏における局所化。}
\VerdierSemanticAnchor{toc-ch2-sec2-1}%
\VerdierContentsRow{.5em}{\textbf{2.1.}\enspace 射の乗法的系。}{111}
\VerdierSemanticAnchor{toc-ch2-sec2-2}%
\VerdierContentsRow{.5em}{\textbf{2.2.}\enspace 局所化圏の構成。}{117}
\VerdierSemanticAnchor{toc-ch2-sec2-3}%
\VerdierContentsRow{.5em}{\textbf{2.3.}\enspace 局所化関手の性質。}{125}

\VerdierSemanticAnchor{toc-ch2-sec3}%
\VerdierContentsSection{3}{三角圏のアーベル化。}
\VerdierSemanticAnchor{toc-ch2-sec3-1}%
\VerdierContentsRow{.5em}{\textbf{3.1.}\enspace 圏 $\mathrm A(\mathcal D)$ の構成。}{135}
\VerdierSemanticAnchor{toc-ch2-sec3-2}%
\VerdierContentsRow{.5em}{\textbf{3.2.}\enspace 圏 $\mathrm A(\mathcal D)$ の性質。}{137}

\VerdierSemanticAnchor{toc-ch2-sec4}%
\VerdierContentsSection{4}{スペクトル対象。}
\VerdierSemanticAnchor{toc-ch2-sec4-1}%
\VerdierContentsRow{.5em}{\textbf{4.1.}\enspace スペクトル対象の定義。}{145}
\VerdierSemanticAnchor{toc-ch2-sec4-2}%
\VerdierContentsRow{.5em}{\textbf{4.2.}\enspace スペクトル系列の機構。}{147}
\VerdierSemanticAnchor{toc-ch2-sec4-3}%
\VerdierContentsRow{.5em}{\textbf{4.3.}\enspace 項 $E\!\begin{bmatrix}n\\-\infty,p,q,+\infty\end{bmatrix}$。}{153}
\VerdierSemanticAnchor{toc-ch2-sec4-4}%
\VerdierContentsRow{.5em}{\textbf{4.4.}\enspace 収束。}{155}

\VerdierSemanticAnchor{toc-ch3}%
\VerdierContentsChapter{第III章：導来圏。}
\VerdierSemanticAnchor{toc-ch3-sec1}%
\VerdierContentsSection{1}{導来圏の定義。}
\VerdierSemanticAnchor{toc-ch3-sec1-1}%
\VerdierContentsRow{.5em}{\textbf{1.1.}\enspace 記法と用語。}{157}
\VerdierSemanticAnchor{toc-ch3-sec1-2}%
\VerdierContentsRow{.5em}{\textbf{1.2.}\enspace 導来圏の重要な部分圏。}{159}
\VerdierSemanticAnchor{toc-ch3-sec1-3}%
\VerdierContentsRow{.5em}{\textbf{1.3.}\enspace 標準 $\delta$-関手。}{164}

\VerdierSemanticAnchor{toc-ch3-sec2}%
\VerdierContentsSection{2}{分解。}
\VerdierSemanticAnchor{toc-ch3-sec2-1}%
\VerdierContentsRow{.5em}{\textbf{2.1.}\enspace 初等射。}{169}
\VerdierSemanticAnchor{toc-ch3-sec2-2}%
\VerdierContentsRow{.5em}{\textbf{2.2.}\enspace ある対象集合の考察。}{173}
\VerdierSemanticAnchor{toc-ch3-sec2-3}%
\VerdierContentsRow{.5em}{\textbf{2.3.}\enspace 分解に関するいくつかの補題。}{177}
\VerdierSemanticAnchor{toc-ch3-sec2-4}%
\VerdierContentsRow{.5em}{\textbf{2.4.}\enspace 応用。}{185}

\VerdierSemanticAnchor{toc-ch3-sec3}%
\VerdierContentsSection{3}{$\operatorname{Ext}$ の研究。}
\VerdierSemanticAnchor{toc-ch3-sec3-1}%
\VerdierContentsRow{.5em}{\textbf{3.1.}\enspace 種々の同値な定義。}{191}
\endgroup
